\begin{hyperbolicbox}{Normalization of lattices}
    \small
    \begin{defi}
        Since we only care about $\Lambda$ up to homothety, for a lattice $\group{p,q}$ with $\tau=\frac{p}{q}\in\H$, we can normalize it to a homothetic lattice $\Lambda_{\tau}=\group{\tau,1}=\frac{1}{q}\group{p,q}$. \cite{silvermanAT} 
    \end{defi}
    
    However, there might also be homothetic normalized lattices such as $\group{i,1}=\group{i+1,1}$. We identify them using the following theorem.

    \begin{thm}
        Let $\tau,\tau'\in\H$. $\Lambda_\tau$ and $\Lambda_{\tau'}$ are homothetic if and only if $\tau$ and $\tau'$ are related with a M\"obius map $\tau'=\frac{a\tau+b}{c\tau+d}$ where $\begin{pmatrix}
            a&b\\c&d
        \end{pmatrix}\in\operatorname{PSL}(2,\Z)$.
    \end{thm}
    \begin{proof}
        By definition, $\Lambda_\tau$ and $\Lambda_{\tau'}$ are homothetic if and only if there is $z\in\C\setminus\set{0}$ such that $\Lambda_{\tau'}=z\Lambda_\tau$, which is equivalent to
        $$
        \group{\tau',1}=\group{z\tau,z}
        $$
        which by \hyperref[thm5]{Theorem 5} holds if and only if there is $\begin{pmatrix}
            a&b\\c&d
        \end{pmatrix}\in\operatorname{SL}(2,\Z)$ such that
        $$
        \begin{pmatrix}
            \tau'\\1
        \end{pmatrix}=\begin{pmatrix}
            a&b\\c&d
        \end{pmatrix}\begin{pmatrix}
            z\tau\\z
        \end{pmatrix}\iff \tau'=\frac{a\tau+b}{c\tau+d}
        $$
    \end{proof}

    This theorem points out the relationship between the action of $\operatorname{PSL}(2,\Z)$ on $\H$ and the classification of $\Lambda\subset\C$.
    
    Specifically, if we use $\mathcal{L}$ to denote the set of all lattices in $\C$, and define an equivalence relation $\sim$ on $L$ by $\Lambda\sim\Lambda'$ if and only if $\Lambda$ and $\Lambda'$ are homothetic, then there is a bijection from $\H/\operatorname{PSL}(2,\Z)$ to $\mathcal{L}/\sim$ given by
    $$
        f:\tau\mapsto\Lambda_\tau
    $$
\end{hyperbolicbox}